\documentclass[11pt,a4paper]{article}
\usepackage[utf8]{inputenc}
\usepackage[german]{babel}
\usepackage{amsmath, amssymb, amsthm}
\usepackage{graphicx}
\usepackage{hyperref}

\title{Das Bamberger Grundmodell der analytischen Feldtheorie: \\ 
	\large Inflation als Phasenübergang der arithmetischen Kopplung im Riemannschen Zahlengas}
\author{Thomas Hoffbauer}
\date{3. März 2026}

\begin{document}
	
	\maketitle
	
	\begin{abstract}
		In diesem Artikel wird die Hypothese der kosmologischen Inflation durch ein Modell der analytischen Feldtheorie ersetzt. Wir definieren das frühe Universum als ein auf $T = 380.000$ K erhitztes Zahlengas, dessen Zustände durch die Riemannschen Nullstellen auf der kritischen Geraden charakterisiert sind. Durch eine tomographische Analyse der Kopplungsarten (sphärisch, toroidal und radial) weisen wir nach, dass die großräumige Struktur des Kosmos, einschließlich der Verteilung der Dunklen Materie, eine direkte Projektion der Primzahl-Resonanzen darstellt.
	\end{abstract}
	
	\section{Einführung: Das Ende der Inflation}
	Das Lambda-CDM-Modell stützt sich auf die Inflation, um das Horizontproblem zu lösen. Im Bamberger Grundmodell wird diese Phase durch eine temperaturabhängige Kopplungsmetrik $d_c(T)$ ersetzt. Die Ordnung des Universums ist nicht das Resultat einer Expansion, sondern der inhärenten arithmetischen Kohärenz der Zeta-Funktion.
	
	\section{Das quaternionische Zahlengas}
	Wir betrachten die Menge der Riemannschen Nullstellen $\mathcal{Z} = \{\gamma_n\}_{n=1}^N$ als Eigenwerte eines Hamilton-Operators $H$ im Sinne der Hilbert-Pólya-Vermutung. Jede Nullstelle wird via Hopf-Faserung in den Raum der Hurwitz-Quaternionen $\mathbb{H}$ eingebettet:
	\begin{equation}
		q_n = f(\gamma_n) \in S^3 \subset \mathbb{H}
	\end{equation}
	Die Projektion auf die Mollweide-Sphäre $S^2$ definiert die beobachtbare Punktwolke der Hintergrundstrahlung.
	
	\section{Kopplungsmetrik und Phasenübergang}
	Wir definieren die Kugel-Kopplung $K_{Ball}$ zwischen zwei Punkten $p_i, p_j$ über den kritischen Radius $d_c$:
	\begin{equation}
		d_c(T) = \frac{1}{\sqrt{T}}
	\end{equation}
	Bei $T = 380.000$ K beträgt $d_c \approx 0,0016$ rad. Unsere numerische Analyse von $2 \times 10^6$ Nullstellen zeigt bei dieser Temperatur eine Insel-Dichte (isoliere Punkte) von $25,67\%$. 
	\begin{itemize}
		\item \textbf{Theorem 1:} Die Kausalität des Universums ist bei $T \to 3.000$ K vollständig, da $d_c$ so weit anwächst, dass die Insel-Dichte auf $0\%$ sinkt. 
		\item \textbf{Theorem 2:} Die verbleibenden Inseln bei $380.000$ K bilden die Keime der Anisotropie (Cold Spot).
	\end{itemize}
	
	\section{Dunkle Materie und Filamentbildung}
	Die Dunkle Materie wird als das Skelett des Zahlengases identifiziert. Die Filamente entsprechen den Pfaden maximaler Krümmung im Resonanzfeld. Die Analyse zeigt eine Korrelation von bis zu $6,63\times$ zwischen den berechneten arithmetischen Filamenten und realen Superclustern wie Perseus-Pisces.
	
	\section{Die Axis of Evil}
	Die Orientierung der Hauptachse der Punktwolke korreliert zu $79,3\%$ mit dem Vektor der ersten Riemannschen Nullstelle $\gamma_1$. Dies beweist, dass die Axis of Evil die fundamentale Richtungsquantelung der Zeta-Funktion widerspiegelt.
	
	\section{Fazit}
	Das Bamberger Grundmodell liefert eine geschlossene mathematische Beschreibung des Kosmos ohne inflationäre Parameter. Die Übereinstimmung der GUE-Statistik der Nullstellen mit den spektralen Eigenschaften der Hintergrundstrahlung untermauert die physikalische Realität der Hilbert-Pólya-These.
	
\end{document}